\documentclass[11pt,letterpaper]{article}
\usepackage[T1]{fontenc}
\usepackage{lmodern,amsmath,amssymb,amsthm,microtype}
\usepackage[margin=1in]{geometry}
\usepackage[hidelinks]{hyperref}
\providecommand{\tightlist}{\setlength{\itemsep}{0pt}\setlength{\parskip}{0pt}}
\DeclareUnicodeCharacter{2264}{\ensuremath{\le}}
\DeclareUnicodeCharacter{2265}{\ensuremath{\ge}}
\DeclareUnicodeCharacter{2260}{\ensuremath{\ne}}
\DeclareUnicodeCharacter{2212}{\ensuremath{-}}

\DeclareUnicodeCharacter{2211}{\ensuremath{\sum}}

\usepackage{longtable,booktabs,array}
\title{A permutation-orbit cut for nine separated subset sums below 150}
\author{Aafi Khan Malek\\ \small Independent researcher, Ahmedabad, India\\
\small\href{mailto:aafimalek2032@gmail.com}{aafimalek2032@gmail.com}}
\date{28 November 2026}
\begin{document}
\maketitle
\textbf{Theorem.} For nine sorted positive real weights with unit-separated
subset sums, \(149\le x_9<150\) implies \(x_8\ge131\).
This is a direct exclusion of the cap \(\lfloor x_8\rfloor\le130\).
The proof uses the
\href{https://github.com/Aafimalek/novel-discoveries/tree/ddc090f4b5de263571dcd4e0ed60b98dbfc387e6/nine-separated-subset-sums/verification/rank-cuts-below-150/penultimate-cap-127}{previous rank-run and interval-matching argument},
with cap-specific exact checks and an orbit argument for the last remaining
cell. This does not close the full interval below 150.

The forced rank edges run from 141 through 172, giving slots
\(q=33,\dots,64\). The base unsigned magnitudes at these slots obey
\(\alpha_q=u_{a-q+1}\in[149-2q,151-2q)\), for common \(64\le a\le74\)
and \(b=148-a\). The bins are:

\begin{longtable}[]{@{}lr@{}}
\toprule\noalign{}
Half-open range & Capacity \\
\midrule\noalign{}
\endhead
\bottomrule\noalign{}
\endlastfoot
{[}1,21) & \(a-64\) \\
each {[}149−2q,151−2q), q=64 down to 33 & 1 \\
{[}85,215) & 32 \\
{[}278,321) & \(148-2a\) \\
{[}383,10000) & \(a-20\) \\
\end{longtable}

The two wide gaps at ranks \(a,b\) have width at least
\(\alpha_{33}-\alpha_{64}+4\ge66\), by the plus-two endpoints. As
before, flips at any weight below 33 preserve membership at both cuts.
Their orbit sizes divide both \(a\) and \(b\), giving \(x_3\ge33\). At
the cap vector \(y=(124,\dots,131)\), balanced forms have magnitude at
most 16 and unbalanced forms at least 240. If any prescribed form is
unbalanced then the total deficit \(R=\sum_{i\le8}(y_i-x_i)>155\)
because every prescribed form is below 85. Otherwise
\(\alpha_{33}\ge83\) is balanced and \(R\ge67\). Thus \(R\ge67\) and the
third deficit bound \(r_3\le126-33=93\) are used.

\section{Exact matching and one final permutation
cut}\label{exact-matching-and-one-final-permutation-cut}

\begin{itemize}
\tightlist
\item
  The initial 7,626 unit cells leave 6,718 after tail/slot matching.
\item
  Splitting by the eleven counts \(a\) gives 73,898 profiles; all-sign
  interval matching leaves 3,239.
\item
  Divisibility removes 36 of those profiles without refinement. The rest
  give 115,315 third-deficit unit cells. The moment/tail bounds remove
  670 and all-sign capacity matching removes 114,644. The unique
  remaining cell has \(a=64\), \(r_1\in[49,50]\), \(r_2,r_3\in[24,25]\).
\end{itemize}

For this cell, the first two adjacent gaps between weights are at most
27 and 2 respectively. The others are at most 26 since the deficits are
nonnegative nonincreasing and \(r_3\le25\). All seven adjacent gaps are
therefore strictly below 33. Swapping signs of adjacent weights changes
the absolute magnitude by at most twice their weight gap, which is
strictly below 66. It therefore cannot change membership at either wide
magnitude gap. Since adjacent swaps generate all coordinate
permutations, both subsets of signed forms counted by \(a\) and \(b\)
are unions of the permutation orbits on the 128 complementary sign
pairs.

These five orbits are indexed by the smaller of the number of plus and
minus signs: their sizes are \(1,8,28,56,35\) from the binomial
coefficients, halving the four-plus orbit at complementation. The union
for \(a\) must be included in the union for \(b\). Assigning each orbit
to the first union, the second union only, or neither gives \(3^5\)
finite choices. None gives \((a,b)=(64,84)\). This excludes the last
residual cell and completes the cap.


\section{Reproducibility}
Run \texttt{python3 finish130.py} in the proof directory with standard Python.
The script performs an exact integer and rational replay, then checks all
\(3^5\) orbit assignments. The result is not verified in Lean.
\end{document}
